\chapter{风险分析与应对措施}

\section{技术风险与应对}

\textbf{风险描述}：田间真实场景复杂（光照、遮挡、复杂背景、虫态差异），实验室训练的模型在实际田间识别精度可能下降；数据采集与标注成本高、周期长。

\textbf{应对措施}：采用「公开数据集预训练 + 自有田间数据微调」策略，优先选用 PlantDoc、IP102 等田间真实数据作为迁移源；引入 YOLOv8-nano 目标检测与 CBAM 注意力机制提升复杂场景鲁棒性；以联邦学习低成本汇聚多区域数据提升泛化性；设定「识别置信度低于阈值则转人工/专家复核」的兜底机制，避免误诊误导。

\section{市场接受度风险与应对}

\textbf{风险描述}：小农户老龄化、数字化接受度低，可能面临「产品好用但用不起来」的冷启动困难。

\textbf{应对措施}：坚持「零学习成本」极简设计，30 秒内上手；借力农资店店主、基层农技员等农户信任的「意见领袖」引导首单体验；以「能挽回损失」的真实案例做口碑传播；初期免费、先以政府示范项目背书降低信任门槛。

\section{运营落地风险与应对}

\textbf{风险描述}：县域下沉市场渠道分散、推广半径长，规模化复制可能面临执行与管理难度。

\textbf{应对措施}：以「样板县——标准化 SOP——区域合伙人」模式推进，形成可复制的县域打法；与农资连锁、供销社等成熟渠道共建利益机制；总部建设数据化运营体系，实时监控获客成本、留存与活跃等核心指标，及时纠偏。

\section{政策与合规风险与应对}

\textbf{风险描述}：农药推荐涉及农资监管与广告合规，识别误判可能引发用药风险与责任纠纷；数据采集涉及个人信息与农业数据安全合规。

\textbf{应对措施}：防治方案严格依据农业农村部农药登记目录与植保技术规范，仅推荐合规登记用药并附免责声明；建立专家审校与方案动态更新机制；数据采集遵循《数据安全法》《个人信息保护法》，端侧优先、最小必要、分级脱敏；购买产品责任险以对冲潜在风险。
